\frontmatterpage
\unchapter{Abstract}

The rapid advancement of mobile technology has transformed how people interact with hospitality services worldwide. In Uzbekistan, traditional teahouses known as ``choyxonas'' represent centuries of cultural heritage, yet their booking processes remain largely manual, relying on phone calls and walk-in visits. This gap between digital expectations and available solutions motivates the development of a specialized mobile application for the Uzbek hospitality market.

The Bachelor Paper consists of four chapters. Chapter~1 reviews the theoretical foundations of mobile application development, analyzes existing restaurant booking platforms such as OpenTable, TheFork, and Resy, and examines the Flutter framework and Firebase backend services that form the technological basis of the project. Chapter~2 presents the system architecture and technical design of the CHOYXONA.UZ application, including database structure, user interface design principles, real-time synchronization mechanisms, and geolocation services integration. Chapter~3 details the functional capabilities of the implemented application across three user roles: customers, establishment administrators, and platform super administrators, covering authentication, search, booking, analytics, and notification subsystems. Chapter~4 addresses testing strategies including unit, widget, and integration testing, the multi-platform deployment process across Google Play, App Store, and thirteen alternative marketplaces, and outlines the future development roadmap.

The methodological basis of the study includes analysis of scientific literature, comparative analysis of existing platforms, iterative software development methodology, and practical implementation with subsequent testing and evaluation. The applied methods include comparative analysis of cross-platform frameworks, prototyping, performance benchmarking, and user acceptance testing.

The main contribution of this work is the design and implementation of a fully functional cross-platform mobile application that bridges traditional Uzbek hospitality culture with modern digital convenience. Performance testing confirmed application responsiveness with 2.3-second cold start time, 60~fps scrolling, and sub-150~ms database query latency. The application was successfully deployed to 15~distribution channels across 248~countries.

\thesisscopeEN

\textbf{Keywords:} mobile application, Flutter, Firebase, restaurant booking, automation, cross-platform development, Uzbekistan.
